% apendices/indice-kanjis-n4-n3.tex
\chapter{\ruby{漢字}{かんじ}\ruby{索引}{さくいん} N4--N3 — Índice de Kanjis}

\begin{tcolorbox}[vocabbox, title={Kanjis N4 — Unidades 01--20 (orden de aparición)}]
\begin{tabularx}{\linewidth}{clXX}
\toprule
\textbf{Kanji} & \textbf{Lectura} & \textbf{Significado} & \textbf{Unidad} \\ \midrule
\ruby{復}{ふく} & ふく & repetir & U01 \\
\ruby{習}{しゅう} & なら(う) & practicar & U01 \\
\ruby{終}{しゅう} & お(わる) & terminar & U02 \\
\ruby{済}{さい} & す(む) & acabar & U02 \\
\ruby{備}{び} & そな(える) & preparar & U03 \\
\ruby{余}{よ} & あま(る) & sobrar & U03 \\
\ruby{済}{さい} & す(む) & terminar & U04 \\
\ruby{変}{へん} & か(わる) & cambiar & U04 \\
\ruby{努}{ど} & つと(める) & esforzarse & U05 \\
\ruby{向}{こう} & む(く) & dirigirse & U05 \\
\ruby{守}{しゅ} & まも(る) & proteger & U06 \\
\ruby{害}{がい} & — & daño & U06 \\
\ruby{降}{こう} & ふ(る) & llover/bajar & U07 \\
\ruby{被}{ひ} & こうむ(る) & sufrir & U07 \\
\ruby{強}{きょう} & つよ(い) & fuerte & U08 \\
\ruby{制}{せい} & — & control & U08 \\
\ruby{我}{が} & われ & yo/uno mismo & U09 \\
\ruby{慢}{まん} & — & tolerancia & U09 \\
\ruby{的}{てき} & まと & objetivo & U10 \\
\ruby{標}{ひょう} & — & marca & U10 \\
\ruby{様}{よう} & さま & forma/aspecto & U11 \\
\ruby{似}{に} & に(る) & parecerse & U11 \\
\ruby{事}{じ} & こと & asunto & U12 \\
\ruby{物}{もの} & もの & cosa & U12 \\
\ruby{証}{しょう} & — & evidencia & U13 \\
\ruby{拠}{きょ} & — & base & U13 \\
\ruby{期}{き} & — & período & U14 \\
\ruby{待}{たい} & ま(つ) & esperar & U14 \\
\ruby{訳}{やく} & わけ & razón/traducción & U15 \\
\ruby{訳}{やく} & — & — & U15 \\
\ruby{限}{げん} & かぎ(る) & limitar & U16 \\
\ruby{除}{じょ} & のぞ(く) & excluir & U16 \\
\ruby{過}{か} & す(ぎる) & exceder & U17 \\
\ruby{度}{ど} & — & grado & U17 \\
\ruby{易}{い} & やさ(しい) & fácil & U18 \\
\ruby{困}{こん} & こま(る) & estar en apuros & U18 \\
\ruby{変}{へん} & か(わる) & cambiar & U19 \\
\ruby{増}{ぞう} & ふ(える) & aumentar & U19 \\
\ruby{原}{げん} & もと & origen & U20 \\
\ruby{因}{いん} & よ(る) & causa & U20 \\
\bottomrule
\end{tabularx}
\end{tcolorbox}

\begin{tcolorbox}[vocabbox, title={Kanjis N3 — Unidades 21--35 (orden de aparición)}]
\begin{tabularx}{\linewidth}{clXX}
\toprule
\textbf{Kanji} & \textbf{Lectura} & \textbf{Significado} & \textbf{Unidad} \\ \midrule
\ruby{当}{とう} & あ(たる) & acertar & U21 \\
\ruby{卒}{そつ} & — & graduarse & U21 \\
\ruby{規}{き} & — & regla & U22 \\
\ruby{成}{せい} & な(る) & lograr & U22 \\
\ruby{相}{そう} & あい & mutuo & U23 \\
\ruby{提}{てい} & さ(げる) & proponer & U23 \\
\ruby{押}{おう} & お(す) & empujar & U24 \\
\ruby{引}{いん} & ひ(く) & jalar & U24 \\
\ruby{思}{し} & おも(う) & pensar & U25 \\
\ruby{考}{こう} & かんが(える) & reflexionar & U25 \\
\ruby{感}{かん} & — & sentir & U25 \\
\ruby{噂}{そん} & うわさ & rumor & U26 \\
\ruby{伝}{でん} & つた(える) & transmitir & U26 \\
\ruby{報}{ほう} & むく(いる) & reportar & U26 \\
\ruby{逆}{ぎゃく} & さか & contrario & U27 \\
\ruby{接}{せつ} & つ(ぐ) & conectar & U27 \\
\ruby{結}{けつ} & むす(ぶ) & anudar & U28 \\
\ruby{影}{えい} & かげ & sombra & U28 \\
\ruby{尊}{そん} & たっと(い) & venerable & U29 \\
\ruby{敬}{けい} & うやま(う) & respetar & U29 \\
\ruby{謙}{けん} & — & modesto & U30 \\
\ruby{拝}{はい} & おが(む) & venerar & U30 \\
\ruby{文}{ぶん} & ふみ & escritura & U31 \\
\ruby{論}{ろん} & — & tesis & U31 \\
\ruby{縮}{しゅく} & ちぢ(む) & contraer & U32 \\
\ruby{若}{じゃく} & わか(い) & joven & U32 \\
\ruby{冊}{さつ} & — & libro & U33 \\
\ruby{枚}{まい} & — & hoja plana & U33 \\
\ruby{読}{どく} & よ(む) & leer & U34 \\
\ruby{段}{だん} & — & párrafo & U34 \\
\ruby{総}{そう} & — & total & U35 \\
\ruby{継}{けい} & つ(ぐ) & continuar & U35 \\
\bottomrule
\end{tabularx}
\end{tcolorbox}

\begin{tcolorbox}[vocabbox, title={\ruby{漢字}{かんじ}の\ruby{部首}{ぶしゅ}別\ruby{グループ}{グループ} — Por radicales clave}]
\begin{tabularx}{\linewidth}{lX}
\toprule
\textbf{Radical} & \textbf{Kanjis del libro} \\ \midrule
\ruby{言}{げん} (言葉) & \ruby{話}{はなし}・\ruby{読}{よむ}・\ruby{語}{かたる}・\ruby{訓}{くん}・\ruby{議}{ぎ}・\ruby{記}{き}・\ruby{詩}{し} \\
\ruby{心}{しん} (mente) & \ruby{思}{おもう}・\ruby{感}{かんじる}・\ruby{悩}{なやむ}・\ruby{忘}{わすれる}・\ruby{慢}{まん} \\
\ruby{手}{しゅ} (mano) & \ruby{持}{もつ}・\ruby{押}{おす}・\ruby{引}{ひく}・\ruby{拝}{はい} \\
\ruby{水}{すい} (agua) & \ruby{海}{うみ}・\ruby{河}{かわ}・\ruby{波}{なみ}・\ruby{泳}{およぐ} \\
\ruby{木}{もく} (árbol) & \ruby{木}{き}・\ruby{林}{はやし}・\ruby{森}{もり}・\ruby{橋}{はし} \\
\ruby{人}{にん} (persona) & \ruby{休}{やすむ}・\ruby{仕}{し}・\ruby{働}{はたらく}・\ruby{仮}{かり} \\
\bottomrule
\end{tabularx}
\end{tcolorbox}
